%%% Document Setup

\newcommand{\templateVersion}{Version 3.0.0}

\usepackage[a4paper, top=25mm, bottom=25mm, left=12.5mm, right=12.5mm, headheight=15pt, headsep=8mm, footskip=10mm]{geometry} % Sets up the page margins.
\usepackage{fancyhdr} % Creates more customisable headers and footers.

\fancypagestyle{plain}{
    \fancyhf{}

    \fancyhead[L]{\titlePageTitle}
    \fancyhead[R]{\leftmark}

    \fancyfoot[L]{Written by \titlePageAuthor. \\ 
    \href{https://github.com/MatthewEmer/LaTeX-File-Template}{\LaTeX\ File Template \templateVersion. \copyright\ \the\year\ Matthew Emerson.}} % Do not edit.
    \fancyfoot[R]{Page \thepage \, of\, {\hypersetup{hidelinks}\pageref{mainDocumentEnd}\relax}}

    \renewcommand{\headrulewidth}{0.4pt}
    \renewcommand{\footrulewidth}{0.4pt}
}

%%%



%%% Title Formatting

\usepackage{titlesec} % Introduces tools for formatting the inbuilt LaTeX section titles.

\titleformat{\chapter}[display]{\normalfont\bfseries}{\large\chaptertitlename\ \thechapter}{6pt}{\Huge}
\titlespacing*{\chapter}{0pt}{-10pt}{20pt}

%%%



%%% Date Formatting

\usepackage{datetime} % Allows greater control over the format of dates.

\newdateformat{monthyeardate} % The date format for the title page.
{%
    \monthname[\THEMONTH] \THEYEAR
}
\newdateformat{yeardate} % The date format for the copyright statement.
{%
    \THEYEAR
}

%%%



%%% Colour Formatting

\usepackage{xcolor} % Introduces tools for colouring text.
\usepackage{hyperref} % Allows for greater link customisation.

\hypersetup{
    colorlinks=true,
    linkcolor=InternalLinkColour,
    citecolor=CitationLinkColour,
    urlcolor=ExternalLinkColour
}

%%%



%%% Citation Formatting

\usepackage[
    backend=biber,
    style=authoryear,
    sorting=none,
    natbib=true
]{biblatex}

\setlength{\bibitemsep}{0.8\baselineskip}
\setlength{\bibhang}{0pt}

\DeclareCiteCommand{\plaincitation}{\usebibmacro{prenote}}
{%
    \printtext[bibhyperref]{%
        \printnames{labelname}%
        \setunit{\space(}%
        \printfield{labelyear})%
    }%
}{\multicitedelim}{\usebibmacro{postnote}}

\DeclareCiteCommand{\bracketcitation}{\usebibmacro{prenote}}
{%
    \printtext[bibhyperref]{%
        (\printnames{labelname}%
        \setunit{\addcomma\space}%
        \printfield{labelyear})%
    }%
}{\multicitedelim}{\usebibmacro{postnote}}

%%%



%%% List of Tables/Figures Formatting

\makeatletter

% List of Tables: no new page, title, or extra heading spacing
\renewcommand{\listoftables}{%
  \@starttoc{lot}%
}

% List of Figures: no new page, title, or extra heading spacing
\renewcommand{\listoffigures}{%
  \@starttoc{lof}%
}

\makeatother

%%%



%%% Additional Packages

\usepackage{float} % Gives additional controls for placing elements in the document.
\usepackage{graphicx} % Additional support for images.
\usepackage{makecell} % Introduces multi-line cells for tables.

\usepackage{dirtytalk} % Fixes speech and quotation marks.
\usepackage{enumitem} % Introduces additional formatting for lists.

\usepackage{amsmath} % Adds in more mathematics functions.
\usepackage{amssymb} % Introduces a wider range of mathematics symbols.

\usepackage{algpseudocode} % Introduces pseudocode snippets.
\usepackage{algorithm} % Allows for algorithms to be inserted.
\usepackage{minted} % Introduces tools for inserting code snippets into the document.

%%%



%%% Text Formatting

\newcommand{\separate}[1]{\vspace{#1 px}\noindent} % #1 Vertical separation in pixels.

\newcommand{\colouredtext}[2]{\color{#1}#2\color{black}} % #1 Text colour, #2 Text to be coloured. 
\newcommand{\comment}[1]{\colouredtext{red}{#1}} % #1 Message. 

%%%



%%% Mathematics 

\newcommand{\N}{\mathbb{N}}
\newcommand{\Z}{\mathbb{Z}}
\newcommand{\Q}{\mathbb{Q}}
\newcommand{\R}{\mathbb{R}}
\newcommand{\I}{\mathbb{I}}
\newcommand{\C}{\mathbb{C}}

\newcommand{\dotproduct}[2]{\langle\mathbf{#1},\mathbf{#2}\rangle} % #1 Left vector, #2 Right vector.
\newcommand{\jacobian}[1]{\mathbf{J}_{#1}} % #1 The function used to create the Jacobian.
\newcommand{\hessian}[1]{\mathbf{H}_{#1}} % #1 The function used to create the Hessian.
\newcommand{\lagrangian}{\mathcal{L}}

%%%



%%% Custom Elements

\newcommand{\newchapter}[1]{\chapter{#1}\label{chapter:#1}} % #1 Chapter name.
\newcommand{\newsection}[1]{\section{#1}\label{section:#1}} % #1 Section name.
\newcommand{\newsubsection}[1]{\subsection{#1}\label{subsection:#1}} % #1 Subsection name.
\newcommand{\newsubsubsection}[1]{\subsubsection{#1}\label{subsubsection:#1}} % #1 Subsubsection name.

\newcommand{\insertfigure}[3]{ % #1 Caption, #2 Image source, #3 Image width.
    \begin{figure}[H]
        \centering
        \includegraphics[width=#3\linewidth]{#2}
        \caption{#1}\label{figure:#1}
    \end{figure}
}

\newcommand{\currenttablecaption}{}
\counterwithout{table}{chapter}
\newenvironment{newtable}[2]{ % #1 Caption, #2 Column layout.
    \renewcommand{\currenttablecaption}{#1}
    \begin{table}[H]
        \centering
        \begin{tabular}{#2}
}{%
        \end{tabular}
        \caption{\currenttablecaption}\label{table:\currenttablecaption}
    \end{table}
}

\newenvironment{newpseudocode}[4]{ % #1 Caption, #2 Inputs, #3 Outputs.
    \begin{algorithm}[H]
        \caption{#1}\label{pseudocode:#1}
        \begin{algorithmic}
            \State \textbf{Inputs}: #2
            \State \textbf{Outputs}: #3
            \separate{5}
}{%
        \end{algorithmic}
    \end{algorithm}
}

\newenvironment{newsnippet}[1]{ % #1 Caption.
    \refstepcounter{algorithm}\label{snippet:#1}
    \par\addvspace{0.8\baselineskip}
    \hrule

    \separate{2}\textbf{Algorithm \thealgorithm} #1 

    \separate{2}\hrule
}{
    \hrule\par\addvspace{0.8\baselineskip}
}

%%%



%%% Custom Academic Environments

\usepackage{amsthm}
\usepackage[nameinlink, noabbrev]{cleveref}

\newtheoremstyle{academic}
  {6pt}       % Space above
  {6pt}       % Space below
  {\normalfont\upshape} % Body font
  {}          % Indentation
  {\bfseries} % Heading font
  {:}         % Punctuation after heading
  {0.5em}     % Space after heading
  {}          % Heading specification
\theoremstyle{academic}

\newtheorem{definitioninner}{Definition}[subsection]
\newenvironment{definition}[1]{\begin{definitioninner}\label{example:#1}}{\end{definitioninner}}

\newtheorem{exampleinner}{Example}[subsection]
\newenvironment{example}[1]{\begin{exampleinner}\label{example:#1}}{\end{exampleinner}}

\newtheorem{conjectureinner}{Conjecture}[subsection]
\newenvironment{conjecture}[1]{\begin{conjectureinner}\label{conjecture:#1}}{\end{conjectureinner}}

\newtheorem{theoreminner}{Theorem}[subsection]
\newenvironment{theorem}[1]{\begin{theoreminner}\label{theorem:#1}}{\end{theoreminner}}

\newtheorem{lemmainner}{Lemma}[subsection]
\newenvironment{lemma}[1]{\begin{lemmainner}\label{lemma:#1}}{\end{lemmainner}}

\newtheorem{propositioninner}{Proposition}[subsection]
\newenvironment{proposition}[1]{\begin{propositioninner}\label{proposition:#1}}{\end{propositioninner}}

\newcommand{\elemref}[1]{\cref{#1}}